\chapter{Einführung}
\label{ch:Einfuehrung}
\todo[inline]{Einführung schreiben}


\section{Problemstellung und Motivation}
\label{sec:Problemstellung_Motivation}
\todo[inline]{Problemstellung und Motivation überarbeiten}
Im Masterstudiengang \enquote{Robotik und Künstliche Intelligenz in der Produktion} der Hochschule Karlsruhe ist ein 
Forschungs- und Entwicklungsprojekt vorgesehen. Diese Arbeit behandelt die Berechnung von Merkmalspunkten aus 
Segmentierungsmasken für Visual Servoing. Als Ergebnis soll ein Konzept zur robusten Featureextraktion aus Segmentierungsmasken erstellt werden, das für bildbasierte 
Regelungsverfahren in Robotersystemen verbessert. Im Fokus stehen Robustheit und Dynamik, da herkömmliche visuelle Steuerungen 
oft durch schwache Merkmalsqualität limitiert sind.

Die visuelle Regelung gewinnt in der industriellen Robotik zunehmend an Bedeutung. Verfahren des Visual Servoing nutzen 
Kamerabilder zur direkten Bewegungssteuerung und ermöglichen flexible Automatisierung. Die Qualität hängt jedoch stark von 
den verwendeten Bildmerkmalen ab. Klassische Detektoren wie \ac{SIFT} oder \ac{ORB} sind effizient, reagieren aber empfindlich 
auf Skalierungs-, Rotations- und Beleuchtungsänderungen, besonders in komplexen Szenen~\cite{lowe2004distinctive, rublee2011orb}. 
Neuere Arbeiten – etwa \textcite{potje2024xfeat} – stellen mit Accelerated Features (\ac{xFeat}) eine hardwareunabhängige, 
schnelle und präzise Methode vor, die sich auch für mobile Systeme eignet. Gleichzeitig zeigen Studien, dass viele neuronale 
Netze einem \emph{Texture Bias} unterliegen: Objekte werden anhand von Oberflächenmustern statt Formen erkannt, was sie 
anfällig für Veränderungen wie Rauschen oder Cartoonisierung macht~\cite{potje2024xfeat}. Menschen hingegen orientieren sich 
stärker an Gestalt und Kontur (\emph{Shape Bias}). Um den Texture Bias zu reduzieren, können Datenaugmentationsmethoden 
eingesetzt werden, die Netzwerke gezielt auf Formmerkmale trainieren~\cite{potje2024xfeat}.

Vor diesem Hintergrund bieten Segmentierungsmasken neue Möglichkeiten. Moderne Netze wie Mask~\ac{R-CNN} erkennen Objekte 
nicht nur, sondern grenzen sie auch präzise von ihrer Umgebung ab~\cite{he2017maskrcnn}. Diese zusätzlichen semantischen 
Informationen bilden eine solide Grundlage für die Merkmalsextraktion. In Kombination mit effizienten Architekturen 
wie \ac{xFeat} und robustheitssteigernden Trainingsmethoden ergibt sich ein vielversprechender Ansatz für zuverlässigere 
und schnellere visuelle Regelungsverfahren.


\section{Ziel der Arbeit}
\label{sec:Zielstellung}
Das Ziel der Arbeit ist die Entwicklung eines Konzepts zur robusten Feature-Extraktion aus Segmentierungsmasken 
für Visual Servoing. Die Qualität der Merkmale soll so verbessert werden, dass bildbasierte Regelungs- und Greifverfahren 
in Robotersystemen zuverlässiger und dynamischer arbeiten können. Während klassische Feature-Detektoren in komplexen 
Szenen häufig versagen, zeigen neuere Verfahren wie \ac{xFeat} zwar eine höhere Effizienz, sind jedoch stark vom 
\emph{Texture Bias} neuronaler Netze geprägt. Um dieses Problem zu adressieren, untersucht die Arbeit, wie 
Segmentierungsmasken und generative Verfahren kombiniert werden können, um den \emph{Shape Bias} zu fördern und 
formbasierte Merkmale gezielt zu verstärken. Auf diese Weise soll eine Grundlage für robustere und stabilere 
visuelle Regelungsverfahren geschaffen werden.



\section{Aufbau der Arbeit}
\label{sec:Aufbau der Arbeit}
Die Arbeit ist wie folgt strukturiert: Kapitel~\ref{ch:Einfuehrung} führt in die Problemstellung, Motivation und 
Zielsetzung ein. In Kapitel~\ref{ch:Grundlagen} werden die theoretischen Grundlagen und der Stand der Technik 
dargestellt, mit Fokus auf Encoder-Decoder-Architekturen, generative Modelle sowie geometrische und semantische 
Transformationstechniken. Kapitel~\ref{ch:Methodik} beschreibt die Entwicklung der Methodik, einschließlich 
der Analyse des ICCV-Papers \cite{He_2023_ICCV}, der identifizierten Herausforderungen und der abgeleiteten 
Lösungsansätze. Darauf aufbauend wird in Kapitel~\ref{ch:Konzept} das finale Konzept vorgestellt, das die 
Online-Datenaugmentierung, die differenzierbare Kantenextraktion und die angepasste Netzarchitektur umfasst. 
Kapitel~\ref{ch:Ergebnisse} präsentiert die Evaluation der entwickelten Ansätze auf den Datensätzen ImageNet256 
und MegaDepth. Abschließend werden in Kapitel~\ref{ch:Diskussion} die Ergebnisse diskutiert, Limitationen 
aufgezeigt und weiterführende Konzepte für zukünftige Arbeiten vorgestellt.




%% End of File %%